\chapter*{Conclusions and Future Work}
\label{cap:conclusions}
The work developed has made it possible to achieve the main objective set out at the beginning of the report: 
to build a functional prototype to support the analysis of preclinical biomedical images using computer vision and artificial intelligence techniques.
The resulting tool makes it possible to automatically segment the brain and cerebral ischaemia regions in mouse magnetic resonance imaging images, calculate 
associated volumetric measurements, and incorporate explainability mechanisms to analyse the behaviour of the models.

Based on the results obtained, it can be concluded that the proposed approach is suitable as a first automated solution within a supervised workflow.
The system does not replace the judgement of the expert researcher, but it does reduce part of the manual workload associated with the delimitation of regions of interest and provides 
a quantitative output that can be reviewed, corrected, and used as support in biomedical studies.

The brain segmentation model has shown particularly robust behaviour. The results achieved, with a mean 3D Dice coefficient of 0.973 and a mean 3D IoU 
of 0.947, indicate a high correspondence between the automatically generated masks and the reference manual segmentations. In addition, the low volumetric error obtained 
confirms that the model not only adequately reproduces the shape of the brain, but also accurately preserves its total volume. This stability makes brain segmentation 
a reliable basis for the later stages of the pipeline.

Cerebral ischaemia segmentation has presented greater difficulty, as expected given the nature of the problem. Ischaemic lesions show diffuse borders, 
intensity variations, and transition areas whose delimitation partly depends on the operator's judgement. Even so, the model has obtained coherent results, with a mean 3D Dice 
of 0.820 and a mean 3D IoU of 0.717. Although these values are lower than those obtained in brain segmentation, they are reasonable for a task with greater anatomical 
and visual ambiguity. The model makes it possible to locate the main pathological regions and offers useful volumetric estimates as a first automated approximation.

From the point of view of quantification, the system fulfils a relevant function, since it transforms the generated masks into interpretable volumetric measurements.
This output is especially important in the context of the work, since the interest of the end user is not limited to visualising a segmentation, but also to 
obtaining numerical values that make it possible to compare cases, study the evolution of a lesion, or analyse differences between subjects.

The incorporation of Explainable Artificial Intelligence techniques has provided an additional layer of analysis. The maps generated using Grad-CAM, Grad-CAM++, and 
Integrated Gradients have made it possible to study which regions of the image influence the predictions of the models. In correctly segmented cases, the explanations 
obtained are concentrated in areas consistent with the segmented structures, which reinforces the interpretation that the models do not rely solely on artefacts or irrelevant 
external regions. Nevertheless, these techniques should be understood as a complementary inspection tool, not as an independent clinical validation.

Overall, the work demonstrates that the use of U-Net architectures with transfer learning is a viable strategy for addressing biomedical segmentation tasks 
in a limited-data context. The decision to use two differentiated models, one for the brain and another for ischaemia, has made it possible to structure the problem clearly 
and to use anatomical segmentation as a prior step before lesion analysis. This organisation improves the coherence of the system and facilitates its interpretation.

Despite these results, the work presents several limitations. The first and most important is the reduced size of the dataset. Although this situation is common 
in biomedical imaging projects, it limits the generalisation capacity of the models and requires the results to be interpreted as an exploratory validation. A more 
robust evaluation would require the incorporation of a larger number of cases, different acquisition protocols, and greater variability in the morphology and intensity of the lesions.

Another relevant limitation is the dependence on the manual annotations used as reference. In the case of ischaemia, the boundaries of the lesion are not always 
clearly defined, so the ground truth contains a certain degree of subjectivity. This variability directly affects metrics such as Dice and IoU, which penalise 
small spatial differences even when these may be acceptable from an expert point of view. Therefore, the results should be interpreted taking into account that the 
manual segmentation represents an operational reference, not an absolute truth.

It should also be considered that the validation has been performed on a limited number of test cases. Although these cases make it possible to study relevant situations, such as 
control, transient ischaemia, and permanent ischaemia, they are not sufficient to state that the system generalises to all possible scenarios. Before its use in broader 
experimental studies, it would be necessary to perform an external validation with more subjects and, preferably, with review by several experts.

\section{Future Work}

As future work, a first line of improvement would consist of expanding the available dataset. The incorporation of new cases would make it possible to train more 
robust models, reduce the risk of overfitting, and evaluate more accurately the behaviour of the system when faced with lesions of different morphology, size, and intensity. It would also 
be advisable to study data augmentation techniques specific to medical imaging, provided that the transformations applied maintain the anatomical coherence of the images.

A second line of work would be to improve the experimental evaluation. In future versions, larger training, validation, and test partitions should be used, 
as well as cross-validation and separate analyses by lesion type, acquisition time, and image quality. In addition, it would be especially useful to compare the predictions of the 
model with annotations made by several experts, in order to quantify interobserver variability and better contextualise the performance of the automatic system.

From a technical point of view, more advanced architectures than the base U-Net used in this work could be explored, such as U-Net++, Attention U-Net, or models with 
attention mechanisms. These architectures could improve the detection of small lesions or lesions with poorly defined borders. However, their use should be evaluated with caution, 
since an increase in complexity does not necessarily guarantee an improvement in contexts with few annotated data.

Another possible improvement would be to study three-dimensional models. The current approach works slice by slice with 2D images, which reduces the computational cost and facilitates the use 
of pretrained models, but does not fully exploit the spatial continuity between consecutive slices. A 3D architecture could improve the volumetric coherence of 
the predictions, although it would require more data, greater computational capacity, and a more complex training strategy.

It would also be advisable to develop a user interface aimed at the biomedical researcher. Although the delivered notebooks make it possible to reproduce the training, 
inference, and evaluation, an application with a graphical interface would facilitate the loading of volumes, the visualisation of slices, the editing of masks, the calculation of volumes, 
and the export of results. This improvement would favour the integration of the tool into a real working environment.

Finally, the methodology developed could be extended to other organs, lesions, or imaging modalities. Although this work has focused on the brain and cerebral ischaemia 
in mice, the proposed pipeline provides a basis that can be adapted to other biomedical segmentation problems, provided that annotated data are available and that 
the preprocessing, training, and validation stages are appropriately adjusted.

In conclusion, this Bachelor's Degree Final Project has made it possible to develop a first functional tool for the automated analysis of preclinical biomedical images.
The results obtained show robust performance in brain segmentation and promising behaviour in ischaemia segmentation, especially taking into 
account the difficulty of the problem and the limitation of the available data. Although the system requires additional validation and improvements before its application in broader 
experimental environments, it constitutes a valid technical basis for moving towards support tools that reduce the manual workload, improve reproducibility, and facilitate the quantification of 
lesions in biomedical research.
